\chapter{Genomic Surveillance Graphs}\label{app:genomic_networks}

\section{Supplementary Methods}\label{sec:app_ch4_supp_methods}

\begin{table}[htbp]
    \centering
    \caption[Grouping of recorded reasons for testing]{
        Mapping of recorded testing reasons to analytical categories.
        The 36 source values were grouped into 11 categories for the cohort summaries in Chapter~\ref{ch:genomic_networks}. The right-hand column lists source values or summarises closely related variants. Records with no testing reason were assigned to the Missing category.
        }\label{tab:app_test_reason_map}
    \begin{thesistablebody}[\thesisdensetablesetup]{@{}P{0.24\linewidth}P{0.68\linewidth}@{}}
    \toprule
    \textbf{Analytical category} & \textbf{Recorded testing-reason values} \\ %chktex 36
    \midrule
    Symptomatic citizen &
    symptomatic-citizen; I have coronavirus symptoms; I live, work or study in a lockdown area with a coronavirus outbreak \\
    \addlinespace
    Symptomatic essential worker &
    symptomatic-essential-worker; I'm an essential worker; scotland-wales-keyworker; wales-keyworker \\
    \addlinespace
    Contact tracing &
    test-for-contact-tracing; test-for-contact-tracing-app; test-for-contact-self-referral; for-symptomatic-household-member; contact with a person who tested positive and contact-tracer/self-referral variants \\
    \addlinespace
    Confirmatory testing &
    confirmatory-positive-test; confirmatory-other-reason; confirmatory-test-unclear; confirmatory-test-borders; told-to-order-repeat-test \\
    \addlinespace
    Isolation scheme &
    self-isolation-support-grant; isolation-testing-home; isolation-testing-facility \\
    \addlinespace
    Clinical testing &
    gp-healthcare-request; antiviral-order; dental-patient-testing; pre-hospital or pre-surgery testing \\
    \addlinespace
    Surveillance and research &
    zoe-symptom-study; contact-testing-study; events-research-programme; serial-testing; ntrg-member \\
    \addlinespace
    Local outbreak &
    local-council-request; attended-outbreak-venue; community-testing; scotland-university; wales-university \\
    \addlinespace
    Travel &
    green-traveller \\
    \addlinespace
    Other &
    other; none; do-not-know; general-cta-referral; personal-assistant; visiting-professional and school/asymptomatic order variants; any other unrecognised non-missing value \\
    \addlinespace
    Missing &
    No testing reason recorded \\
    \bottomrule
    \end{thesistablebody}
\end{table}

\begin{figure}[htbp]
    \centering
    \includegraphics[width=\textwidth]{fig_simd_population_weighting}
    \caption[Validation of population-weighted SIMD quintiles]{\textbf{Validation of the population-weighted SIMD quintile assignment used in Chapter~\ref{ch:genomic_networks}}. (A) population share in each quintile under the stored SIMD labels, an equal-Data-Zone reconstruction, and the population-weighted reconstruction used for analysis. The stored SIMD labels matched the equal-Data-Zone reconstruction, so these profiles overlap; the population-weighted grouping remains close to the 20\% reference line in each quintile. (B) cross-classification of equal-Data-Zone and population-weighted assignments, with cells showing the percentage of the national population in each combination. Off-diagonal cells indicate Data Zones at quintile boundaries that move when cumulative-population cut-points are used. The accompanying Table~\ref{tab:simd_population_weighting} gives the corresponding Data Zone counts, population totals, population shares, and SIMD rank ranges for the equal-Data-Zone and population-weighted groupings.}\label{fig:app_simd_population_weighting}
\end{figure}

\input{assets/scotland/tab_simd_population_weighting}

\input{assets/policy/policy_periods}

\begin{figure}[htbp]
    \centering
    \includegraphics[width=\textwidth]{policy_index_comparison}
    \caption[Comparison of OxCGRT policy indices for Scotland]{\textbf{Comparison of OxCGRT policy indices for Scotland.} (A) Daily Stringency and Containment and Health Index values. (B) Paired daily values with an ordinary linear fit and identity line. The comparison includes 1,036 complete days from 1~March~2020 to 31~December~2022. Both indices describe formally recorded policy rather than adherence or effectiveness \citep{Hale2021}.}\label{fig:app_policy_index_comparison}
\end{figure}

\begin{figure}[htbp]
    \centering
    \includegraphics[width=\textwidth]{fig_parameter_sensitivity}
    \caption[Parameter sensitivity of clustering and sparsification settings]{\textbf{Genomic graph construction parameter sensitivity.} (A) agreement between Leiden partitions at alternative resolution values and the primary analysis resolution, $R=0.3$, measured by adjusted Rand index (ARI)\nomenclature{ARI}{adjusted Rand index}. (B) clustering fragmentation across Leiden resolutions, shown as clusters per 1,000 sequences relative to the primary resolution. In panels A and B, lines show the median across analysis windows and shaded bands show the interquartile range; the vertical dotted line marks $R=0.3$. (C) sparsification trade-off across EpiLink compatibility thresholds, comparing the pooled fraction of retained pairwise rows with the pooled fraction of retained compatibility weight. The highlighted primary threshold, 0.001, retained 33.1\% of pairwise rows and 99.0\% of total compatibility weight.}\label{fig:app_parameter_sensitivity}
\end{figure}

\clearpage
\subsection{Technical details of assortativity estimation}\label{sec:app_ch4_assortativity_methods}

This section gives the estimation details underlying the compatibility-graph assortativity summaries reported in Chapter~\ref{ch:genomic_networks}. The main text retains the core interpretation of the assortativity coefficient; the equations below document the weighted mixing construction, bootstrap uncertainty, lineage and window pooling, variance decomposition, and supplementary topology assortativity diagnostics.

\subsubsection{Weighted categorical mixing}

For graph $g$, let $V_g$ be the sampled infections and let $L_{gm}$ be the retained undirected compatibility edges whose endpoints both had a non-missing label for attribute $m$. Each edge $e=\{i,j\}$ connecting sequences $i$ and $j$ had EpiLink compatibility weight $w_e$. For categories $a,b$, the symmetric weighted mixing count was

\[
	M^{(g,m)}_{ab}
	=
	\sum_{\{i,j\}\in L_{gm}}
	w_{ij}
	\left\{
	\mathbbm{1}(x_i=a, x_j=b)
	+
	\mathbbm{1}(x_i=b, x_j=a)
	\right\},
\]

and the normalised mixing matrix was

\[
	E^{(g,m)}_{ab}
	=
	\frac{M^{(g,m)}_{ab}}
	{\sum_{c}\sum_{d} M^{(g,m)}_{cd}}.
\]

Each undirected edge therefore contributes in both directions, and $\sum_{a,b}E^{(g,m)}_{ab}=1$. The marginal compatibility-weight share for category $a$ was $p^{(g,m)}_a=\sum_b E^{(g,m)}_{ab}$.

\subsubsection{Multiplier bootstrap uncertainty}

Graph-specific uncertainty was estimated using an edge-weight multiplier bootstrap \citep{Efron1979,Praestgaard1993}. For bootstrap replicate $b=1,\ldots,B$, independent exponential multipliers were drawn for each retained edge,

\[
	\xi^{(b)}_e \sim \operatorname{Exponential}(1),
	\qquad
	w^{*(b)}_e = w_e \xi^{(b)}_e,
\]

and $r^{*(b)}_{gm}$ was recomputed from the perturbed weighted mixing matrix. We used $B=500$ replicates. The graph-level standard error was the sample standard deviation of the finite bootstrap estimates,

\[
	\widehat{\operatorname{SE}}(r_{gm})
	=
	\left[
		\frac{1}{K-1}
		\sum_{b=1}^{K}
		\left(r^{*(b)}_{gm}-\bar r^{*}_{gm}\right)^2 %chktex 3
		\right]^{1/2}, %chktex 3
\]

where $K$ is the number of finite bootstrap replicates. Confidence intervals used the percentile interval,

\[
	\left[
		Q_{\alpha/2}\{r^{*}_{gm}\},
		Q_{1-\alpha/2}\{r^{*}_{gm}\}
		\right],
	\qquad \alpha=0.05.
\]

\subsubsection{Random-effects pooling and overlap-aware GLS}

To account for the fact that lineages vary in size and that windows overlap in time, we used a two-stage hierarchical approach: first pooling lineages within windows, then pooling windows across time while correcting for calendar overlap.

For attribute $m$, window $t$, and lineage graph $\ell$, let $y_{\ell tm}=r_{\ell tm}$ be the observed assortativity and $v_{\ell tm}=\widehat{\operatorname{SE}}(r_{\ell tm})^2$ be its variance within window. We fitted a random-effects model: %chktex 3

\[
	y_{\ell tm}=\theta_{tm}+u_{\ell tm}+\varepsilon_{\ell tm},
	\qquad
	u_{\ell tm}\sim N(0,\tau^2_{tm}),
	\qquad
	\varepsilon_{\ell tm}\sim N(0,v_{\ell tm}),
\]

where $\theta_{tm}$ is the true pooled effect for window $t$, $u_{\ell tm}$ is the between-lineage random effect with variance $\tau^2_{tm}$, and $\varepsilon_{\ell tm}$ is the sampling error. We estimated the between-lineage variance $\tau^2_{tm}$ by restricted maximum likelihood (REML)\nomenclature{REML}{restricted maximum likelihood} \citep{Maestrini2026}. The pooled window estimate was

\[
	\hat\theta_{tm}
	=
	\frac{\sum_{\ell} q_{\ell tm} y_{\ell tm}}
	{\sum_{\ell} q_{\ell tm}},
	\qquad
	q_{\ell tm}=\frac{1}{v_{\ell tm}+\hat\tau^2_{tm}}.
\]

Window-level confidence intervals used a modified Hartung-Knapp adjustment with REML-estimated between-lineage variance \citep{Hartung2001}. We floored the Hartung-Knapp variance estimate at the conventional random-effects inverse-variance estimate to avoid anti-conservative intervals when between-lineage dispersion was small. Intervals were bounded to the valid assortativity range $[-1,1]$.

The overall attribute-level mean across rolling windows was estimated with an intercept-only correlated generalised least-squares (GLS) model \citep{Aitken1936}. Let $\hat{\boldsymbol\theta}_m$ be the vector of window-level pooled estimates and $S_m$ its working covariance matrix. Let $\mathbf{1}$ denote a column vector of ones. The GLS mean was

\[
	\hat\mu_m
	=
	\left(\mathbf{1}^{\mathsf T}S_m^{-1}\mathbf{1}\right)^{-1} %chktex 3
	\mathbf{1}^{\mathsf T}S_m^{-1}\hat{\boldsymbol\theta}_m,
	\qquad
	\operatorname{Var}(\hat\mu_m)
	=
	\left(\mathbf{1}^{\mathsf T}S_m^{-1}\mathbf{1}\right)^{-1}. %chktex 3
\]

For windows $t$ and $t'$, the working correlation was proportional to their calendar overlap,

\[
	\rho_{tt'}
	=
	\frac{|I_t\cap I_{t'}|}
	{\sqrt{|I_t|\,|I_{t'}|}},
	\qquad
	S_{tt'}=\rho_{tt'}\,\widehat{\operatorname{SE}}(\hat\theta_{tm})\,
	\widehat{\operatorname{SE}}(\hat\theta_{t'm}).
\]

We added a negligible diagonal jitter to $S_m$ before inversion for numerical stability. Confidence intervals for the GLS mean used the normal approximation.

\subsubsection{Variance decomposition}

For the variance decomposition, we used graph-level assortativity estimates with finite, positive bootstrap standard errors. We fitted weighted least-squares models with inverse-variance weights $1/v_{\ell tm}$, capped at the 90th percentile to limit domination by very small bootstrap standard errors. For each attribute, we compared intercept-only, window-only, lineage-only, and additive window-plus-lineage models. If $\operatorname{RSS}_0$\nomenclature{RSS}{residual sum of squares}, $\operatorname{RSS}_W$, $\operatorname{RSS}_L$, and $\operatorname{RSS}_{W+L}$ denote the weighted residual sums of squares from these models, the reported fractions included

\[
	F_W=\frac{\operatorname{RSS}_0-\operatorname{RSS}_W}{\operatorname{RSS}_0},
	\quad
	F_L=\frac{\operatorname{RSS}_0-\operatorname{RSS}_L}{\operatorname{RSS}_0},
	\quad
	F_{W+L}=\frac{\operatorname{RSS}_0-\operatorname{RSS}_{W+L}}{\operatorname{RSS}_0},
	\quad
	F_{\mathrm{res}}=\frac{\operatorname{RSS}_{W+L}}{\operatorname{RSS}_0}.
\]

Together, these fractions examine whether graph-level assortativity for each attribute varied primarily across rolling sampling windows, across Pango lineages, through their additive combination, or through residual graph-level heterogeneity not captured by either factor.

\subsubsection{Node degree and strength assortativity}

For node $i$, unweighted degree was

\[
	k_i=\sum_{j}\mathbbm{1}\{\{i,j\}\in L_g\},
\]

and compatibility strength was

\[
	s_i=\sum_{j:\{i,j\}\in L_g} w_{ij}.
\]

For a scalar node attribute $z_i$, such as $k_i$ or $s_i$, numeric assortativity across edges with analysis weights $a_{ij}$ was

\[
	\rho_z
	=
	\frac{
		\frac{\sum_{\{i,j\}} a_{ij}z_i z_j}{\sum_{\{i,j\}} a_{ij}}
		-
		\bar z^2
	}{
		\frac{\sum_{\{i,j\}} a_{ij}(z_i^2+z_j^2)}{2\sum_{\{i,j\}} a_{ij}}
		-
		\bar z^2
	},
	\qquad
	\bar z=
	\frac{\sum_{\{i,j\}} a_{ij}(z_i+z_j)}{2\sum_{\{i,j\}} a_{ij}}.
\]

We reported three diagnostics: degree assortativity with $z_i=k_i$ and $a_{ij}=1$; weighted degree assortativity with $z_i=k_i$ and $a_{ij}=w_{ij}$; and strength assortativity with $z_i=s_i$ and $a_{ij}=w_{ij}$. For strength assortativity, multiplier-bootstrap uncertainty was calculated by recomputing node strengths under the perturbed edge weights.

\clearpage

\section{Supplementary Results}\label{sec:app_ch4_supp_results}

\begin{figure}[htbp]
	\centering
	\includegraphics[width=\textwidth]{fig_assortativity_pooled_window}
	\caption[Pooled window compatibility assortativity over time]{\textbf{Pooled window EpiLink compatibility assortativity by attribute.} Random-effects pooled window mean (solid blue) with 95\% confidence interval (blue band), and the median (dashed orange) and interquartile range (orange band) across lineage graphs within each window, for (A) age group, (B) health board, (C) local authority, (D) sex, (E) SIMD quintile, and (F) urban/rural class. The horizontal dashed line marks zero (random mixing); positive values indicate more within-category compatibility weight than expected.
    Pooled summary assortativity by attribute within window-lineage graphs is presented in Section~\ref{sec:ch4_graph_assortativity}.
    }\label{fig:app_assortativity_pooled_window}
\end{figure}

\begin{figure}[htbp]
	\centering
	\includegraphics[width=\textwidth]{fig_assortativity_variance_decomposition}
	\caption[Variance decomposition of compatibility assortativity]{\textbf{Variance decomposition of window-level assortativity across winsorising percentiles.} Fraction of assortativity variance explained by (A) the additive window-plus-lineage model (dashed = adjusted), (B) window given lineage, and (C) lineage given window, as a function of the winsorising percentile applied to the inverse-variance weights. The vertical dashed line marks the reported 90th-percentile cap.}\label{fig:app_assortativity_variance_decomposition}
\end{figure}

\begin{table}[htbp]
    \centering
    \caption[Diagnostic metrics for variance decomposition]{Diagnostic metrics reported alongside the variance decomposition of compatibility assortativity. These quantities describe the weighting and uncertainty of the underlying window-lineage assortativity estimates. They are included to aid interpretation of Table~\ref{tab:assortativity_variance_decomposition}, but are not themselves variance components.}\label{tab:app_assortativity_variance_decomposition_diagnostics}
    \begin{thesistablebody}{@{}P{0.22\textwidth}P{0.40\textwidth}P{0.30\textwidth}@{}}
    \toprule
    \textbf{Metric} & \textbf{Definition} & \textbf{Interpretation} \\
    \midrule
    Weight cap &
    The 90th-percentile threshold applied to the inverse-variance weights for a given attribute. Any weight larger than this value was truncated to the cap before fitting the decomposition models. &
    Limits the influence of a small number of very precise assortativity estimates. Larger values indicate that some estimates had high precision and therefore large inverse-variance weights. \\
    \addlinespace[0.35em]

    Capped weights &
    The number of window-lineage estimates whose inverse-variance weights exceeded the 90th-percentile threshold and were therefore truncated. &
    Shows how many observations were affected by winsorising. With a 90th-percentile cap, this is expected to be approximately 10\% of the rows for each attribute. \\
    \addlinespace[0.35em]

    Median boot SE &
    The median bootstrap standard error of the underlying window-lineage assortativity estimates for a given attribute. &
    Summarises the typical uncertainty of the assortativity estimates. Larger values indicate noisier estimates and therefore less precise evidence for assortative structure. \\
    \addlinespace[0.35em]

    Median CI width &
    The median width of the bootstrap confidence intervals for the underlying window-lineage assortativity estimates. &
    Provides an alternative summary of typical uncertainty. Wider intervals indicate less precise assortativity estimates. \\
    \bottomrule
    \end{thesistablebody}
\end{table}

\begin{figure}[htbp]
    \centering
    \includegraphics[width=\textwidth]{fig_compatibility_topology}
    \caption[Compatibility-graph topology diagnostics]{\textbf{Compatibility-graph topology diagnostics.} (A) node and retained-edge counts; (B) total retained compatibility weight; and (C) degree assortativity with equal edge weights, degree with weighted edges, and strength assortativity with weighted edges (green ribbon denotes 95\% CI for strength assortativity). Shaded backgrounds denote epidemic eras (cf. Figure~\ref{fig:surveillance_timeline}). Graph scale varied widely (median window summaries: 3,663 nodes; 116,499 retained edges; total retained weight 6,216), with node and retained-edge counts peaking in the window centred on 15~March~2022. Equal-weight degree assortativity was weakly negative on average (median -0.111), whereas weighted degree assortativity and strength assortativity were mostly positive (medians 0.499 and 0.206). These summaries describe sampled compatibility-graph structure, not transmission-graph topology.}\label{fig:app_compatibility_topology}
\end{figure}

\begin{figure}[htbp]
    \centering
    \includegraphics[width=\textwidth]{fig_compatibility_topology_correlations}
    \caption[Compatibility-topology correlations with sequencing denominators]{\textbf{Compatibility-topology correlations with sequencing denominators.} Pearson correlations across rolling windows between sequencing denominators and compatibility-topology summaries. Summaries are window-level averages across lineage-specific graphs, weighted by each graph's total retained EpiLink compatibility weight. Node count is included as a consistency check (correlation with sequence count = 1.00) because the diagnostics use the rolling-window node set. Sequencing volume was strongly correlated with graph-scale summaries, including retained edges (0.80), total retained edge weight (0.81), mean degree (0.76), maximum degree (0.89), mean strength (0.64), and maximum strength (0.86). Positive-test counts showed a similar but weaker pattern. Sequencing coverage was weakly correlated with graph-scale summaries (for example, -0.03 with retained edges, 0.00 with total retained edge weight, and 0.08 with mean degree). Assortativity correlations were smaller and mixed in direction, so they should be interpreted separately from graph scale.}\label{fig:app_compatibility_topology_correlations}
\end{figure}
